\section{Preliminaries II: Algebraic Setting and SIS/ISIS}
\label{sec:prelim-algebra}

\paragraph{Notation.}
Let $q\ge 2$ be an integer modulus and let $\varphi$ be a power of two.
We work in the cyclotomic ring $R \coloneqq \mathbb{Z}[X]/(X^\varphi+1)$ and its $q$-adic reduction
$R_q \coloneqq R/qR$. We use $\varphi$ exclusively for the ring degree and $N$ for the PCS/DECS domain size; when legacy formulas write $n$, read $n=\varphi$.
For vectors over $R$ we use boldface (e.g., $\mathbf{a}\in R^m$), and we write
$\|\cdot\|_2,\ \|\cdot\|_\infty$ for coefficient-wise norms on $R$ (extended component-wise to $R^m$).
Equalities “$\equiv \bmod q$” and Hadamard operations act coefficient-wise in the evaluation basis.
For a matrix $\mathbf{A}\in R_q^{k\times m}$ and target $\mathbf{t}\in R_q^k$, denote by
$\Lambda^\mathbf{t}_\perp(\mathbf{A}) \coloneqq \{\mathbf{e}\in R^m : \mathbf{A}\mathbf{e}\equiv \mathbf{t}\ (\bmod\ q)\}$ the corresponding coset of the lattice.


\paragraph{Convention (coefficient‑wise modeling).}
Via the NTT/evaluation basis we identify $R_q$ with $\F_q^N$; the Hadamard product $\Had$ and division $\Hdiv$ act component‑wise in that basis.
Equalities “$\equiv$ mod $q$” are understood coefficient‑wise unless stated otherwise.

\paragraph{Bounded alphabets and signed lifting.}
We fix an integer bound $B\ge 1$ and write
\[
[-B,B] := \{-B,-B{+}1,\ldots,0,\ldots,B\}\subset\mathbb{Z}.
\]
When encoding integers into $\mathbb{F}_q$, we use the \emph{signed lifting} $\langle i\rangle_q$ defined by
$\langle i\rangle_q = i \pmod q$ for $i\ge 0$ and $\langle i\rangle_q = q - ((-i)\pmod q)$ for $i<0$.

\subsection{SIS and ISIS (Ring Setting) and How We Instantiate Them}


\begin{definition}[SIS and ISIS over $R_q$]\label{def:sis-isis}
Fix integers $n,m\ge 1$ and a modulus $q\ge 2$, and work over $R_q$.
\begin{itemize}[leftmargin=1.25em]
  \item \textbf{SIS (homogeneous).} Given $\mathbf{A}\in R_q^{\,n\times m}$ and a bound $\beta>0$, find a nonzero $\mathbf{s}\in R^{m}$ such that $\mathbf{A}\mathbf{s}\equiv \mathbf{0}\pmod q$ and $\|\mathbf{s}\|\le\beta$ (for a prescribed norm, e.g., $\ell_\infty$ or $\ell_2$).
  \item \textbf{ISIS (inhomogeneous).} Given $(\mathbf{A},\mathbf{t})\in R_q^{\,n\times m}\times R_q^{\,n}$ and $\beta>0$, find $\mathbf{s}\in R^{m}$ with $\mathbf{A}\mathbf{s}\equiv \mathbf{t}\pmod q$ and $\|\mathbf{s}\|\le\beta$.
\end{itemize}
Hardness is measured against PPT adversaries; norms may be coefficient-wise $\ell_\infty$ or Euclidean $\ell_2$ as specified.
\end{definition}

\begin{definition}[Lattice trapdoor suite]\label{def:trapdoor}
Let $R_q=R/qR$. A lattice trapdoor suite consists of two PPT algorithms $(\TrapGen,\SampPre)$:
\begin{itemize}[leftmargin=1.25em]
  \item $(\mathbf{A},\mathsf{td}_\mathbf{A})\leftarrow \TrapGen(1^\lambda,1^n,1^m,q)$ outputs $\mathbf{A}\in R_q^{\,n\times m}$ statistically close to uniform together with a trapdoor $\mathsf{td}_\mathbf{A}$.
  \item $\mathbf{u}\leftarrow \SampPre(\mathsf{td}_\mathbf{A},\mathbf{v},s)$, on input $\mathbf{v}\in R_q^{\,n}$ and Gaussian parameter $s>0$, samples a short $\mathbf{u}\in R^{m}$ such that $\mathbf{A}\mathbf{u}\equiv \mathbf{v}\pmod q$.
\end{itemize}
For admissible parameters (as in GPV-type analyses), the output of $\SampPre$ is statistically close to the discrete Gaussian $D_{\Lambda^{\mathbf{v}}_{\perp}(\mathbf{A}),\,s}$. See, e.g., \cite{GPV08} for concrete admissibility conditions.
\end{definition}

\begin{remark}[Trapdoors]\label{rem:trapdoors-bridge}
In what follows we rely on $(\TrapGen,\SampPre)$ to sample short ISIS preimages and to argue statistical properties of signatures.
For a formal abstract statement see \Cref{def:trap-sampler}; our concrete ring/NTRU instantiation appears in \S\ref{sec:gaussian-preimage}.
\end{remark}

See also the abstract interface in Def.~\Cref{def:trap-sampler}.

\paragraph{From SIS/ISIS to our instantiation.}
Our signature core is an ISIS preimage relation with a BBS-style \emph{rational hash}; its output \(\mathbf{t}\) serves as the commitment image.
Public parameters include a trapdoored matrix \(\mathbf{A}\in R_q^{\,n\times m}\) (verification map) and a
public matrix \(\mathbf{B}\in R_q^{\,n\times \ell}\) (hash domain). For a message \(u\) and signing randomness
\(\chi\), the vSIS construction defines a rational function \(h_{u,\chi}\) and sets the target
\[
  \mathbf{t}\ \coloneqq\ h_{u,\chi}(\mathbf{B})\ \in R_q^{\,n}.
\]
A signature is precisely a short preimage \(\mathbf{s}\) such that
\[
  \boxed{\ \mathbf{A}\mathbf{s}\ \equiv\ \mathbf{t}\ \pmod q\ } \qquad\text{with}\quad \|\mathbf{s}\|\ \text{small}.
\]
This is the verification equation we will prove in zero knowledge later on.
Concrete choices of \(h_{u,\chi}\) translate the verification equations into low-degree
polynomial relations over \(R_q\); e.g., in the BBS-style line one has functions of the form
\(h_{u,\chi}(\tilde{\mathbf{b}}_0,\tilde{\mathbf{b}}_1)
  = \big((1,u^\top,x_0^\top)\cdot \tilde{\mathbf{b}}_0\big)\big/(\tilde{\mathbf{b}}_1-x_1)\),
with \(\chi=(x_0,x_1)\), which exactly yields an ISIS target \(\mathbf{t}\) when evaluated at the public
parameter \(\mathbf{B}=(\mathbf{b}_0,\mathbf{b}_1)\). The vSIS paper formalizes this pattern through a
general signature skeleton \(\Sigma_{\mathrm{ISIS_f}}\) where verification checks \(\mathbf{A}\mathbf{s}=h_{u,\chi}(\mathbf{B})\pmod q\)
and \(\mathbf{s}\) is short; our instantiation follows that template. 

\begin{definition}[Rational-hash commitment]
For message $u$ and signing randomness $\chi=(x_0,x_1)$, define
\[
  t\;\coloneqq\;h_{u,\chi}(B)\in R_q^n
\]
as the \emph{commitment} to $u$. The holder computes $t$ locally and reveals it only as required by the issuance/verification flow.\;\emph{Convention.} We call $h_{u,\chi}$ the \emph{rational hash} and $t$ its \emph{commitment image}.
\end{definition}

\paragraph{Sampling/representation conventions (as in \cite{cryptoeprint:2025/356}).}
Inside the proof, each component is restricted coefficient-wise to a bounded alphabet:
\[\nu \in [-B,B]^{kN},\qquad x_0 \in [-B,B]^{k_0N},\qquad x_1 \in [-B,B]^{N}.\]
(When $B{=}1$ this matches the balanced-ternary alphabet $\{-1,0,1\}$.)
We continue to compute $t=h_{u,\chi}(B)$ exactly in the NTT domain and resample $x_1$ until every coordinate of $b_1-x_1$ is invertible; the failure probability per coordinate is $(2B{+}1)/q$ and remains negligible.

\paragraph{Norm convention.}\label{par:norm-conv}
Unless explicitly stated otherwise, validity is checked via $\ell_\infty$ per coefficient (enforced by the chain gadget in \S\ref{sec:norm-gadget}). We choose the global Euclidean budget to conservatively imply the $\ell_2$ bound via
\[
\beta_2^2\ \ge\ T\,n\,B_\infty^2,
\]
where $T$ is the number of signature rows and $B_\infty$ follows from the chosen $(R,D,L)$ in the chain gadget.


\paragraph{Presentation unlinkability.} See the concrete \textsf{Show} in Section~\Cref{sec:encoding-vsis-smallwood}.

\subsection{Valid Preimage for Our ISIS Instantiation}
\label{subsec:valid-preimage-refined}
Let the public key fix \(\mathbf{A}\in R_q^{\,n\times m}\) and the public parameters fix \(\mathbf{B}\in R_q^{\,n\times \ell}\).
For a message \(u\) and randomness \(\chi\), define the vSIS–BBS hash target
\[
  \mathbf{t}\ \coloneqq\ h_{u,\chi}(\mathbf{B})\ \in R_q^{\,n},
\]
and the admissible witness set (for a chosen norm and bound \(\beta\))
\[
  \mathcal{W}_\beta(u,\chi)\ \coloneqq\ \Bigl\{\,\mathbf{s}\in R^{m}\ :\ \mathbf{A}\mathbf{s}\equiv \mathbf{t}\pmod q,\ \ \|\mathbf{s}\|\le \beta\,\Bigr\}.
\]
\begin{definition}[Valid preimage, ISIS–BBS]
A vector \(\mathbf{s}\in R^{m}\) is a \emph{valid preimage} for \((u,\chi)\) if
\(\mathbf{s}\in \mathcal{W}_\beta(u,\chi)\) i.e.,
\[
  \mathbf{A}\mathbf{s}\equiv h_{u,\chi}(\mathbf{B})\pmod q
  \quad\text{and}\quad
  \|\mathbf{s}\|_\infty\le \beta_\infty\ \ \text{(or }\|\mathbf{s}\|_2\le \beta_2\text{)}.
\]
\end{definition}
The bound \(\beta\) is chosen large enough for correctness under discrete-Gaussian preimage sampling
(\S\ref{subsec:gaussian}), yet small enough so that (i) the induced commitment/hash is computationally
binding, and (ii) our ZK constraints (the \(\ell_\infty\) gadget) remain efficient. In the vSIS
framework, \(\mathbf{t}=h_{u,\chi}(\mathbf{B})\) is produced by an algebraic, keyed rational function family,
and signatures are precisely short preimages for these images. 

\noindent In our BBS instantiation, the message subspace $\mathcal{M}_{X,\mathcal{H},B}$ simply enforces invertibility of each coordinate of $(b_1-x_1)$; we implement this by resampling $x_1$ until $(b_1-x_1)$ is slot-wise invertible (Algorithm~\ref{alg:compute-bbs-exact}, step~4).

\subsection{SmallWood Building Blocks (Quick Definitions)}
\label{subsec:smallwood-defs}
We rely on the following commitment/argument components:
\begin{itemize}[leftmargin=1.25em]
  \item \textbf{DECS} (Degree-Enforcing Commitment Scheme): a commitment layer that enforces low-degree structure of committed vectors/polynomials and supports proximity checks.
  \item \textbf{LVCS} (Linear Vector Commitment System): a vector commitment with linear operations and succinct openings; used to batch and route constraints.
  \item \textbf{PCS} (Polynomial Commitment Scheme): commit to a polynomial and prove evaluations/relations at points.
  \item \textbf{PIOP} (Polynomial IOP): an interactive oracle proof where the prover posts low-degree polynomials and the verifier queries them at random points; compiled via a PCS to a succinct (N)IZK.
\end{itemize}
SmallWood composes these layers (DECS\,$\rightarrow$\,LVCS\,$\rightarrow$\,PCS/PIOP) with hash-based commitments to obtain compact arguments and efficient norm gadgets for our use cases.

\subsection{Discrete Gaussian Preimage Sampling}
\label{subsec:gaussian}
\begin{definition}[Gaussian preimage sampler]\label{def:gaussian-sampler}
Given $(\mathbf{A},\mathsf{td}_\mathbf{A})$ and $\mathbf{t}\in R_q^{k}$, a Gaussian preimage sampler is a randomized algorithm
\[
  \mathsf{SamplePreimage}(\mathbf{A},\mathsf{td}_\mathbf{A},\mathbf{t};s)\ \longrightarrow\ \mathbf{e}\in \Lambda^{\mathbf{t}}_{\perp}(\mathbf{A})
\]
whose output distribution is statistically close to $D_{\Lambda^{\mathbf{t}}_{\perp}(\mathbf{A}),\,s}$ for a width $s$ above a scheme-dependent smoothing threshold. Such samplers (e.g., GPV and ring/NTRU refinements) underpin hash-and-sign signatures with overwhelming correctness and strong statistical properties.
\end{definition}

In practice we instantiate $\mathsf{SamplePreimage}$ with a ring–NTRU trapdoor and per-slot widths calibrated as in \S\ref{sec:gaussian-preimage}; our implementation is constant-time in the coefficient domain.

\lena{
\subsection{Anonymous Credentials}
Anonymous credentials allow users to prove possession of valid credentials without revealing their identity or enabling linking across presentations. A foundational primitive for such systems is a Signature with Efficient Protocols~(SEP)~\cite{SCN:CamLys02}. A SEP is a signature scheme that supports efficient evaluation and ZKPs of possession of a valid signature in a two-party setting. SEPs allow the creation of a digital credential, where a user can prove they possess a valid credential without revealing which specific credential or any identifying information. 

A basic anonymous credential supports only a single presentation: after one use, subsequent presentations become linkable to the first since the same proof would be shown again. To achieve \emph{multi-show credentials} with unlinkable presentations, the construction must incorporate a mechanism to ensure each presentation appears independent while preventing unlimited reuse.

Recent IETF standardization efforts address multi-show credentials through Anonymous Rate-Limited Credentials~(ARC)~\cite{draft-yun-privacypass-crypto-arc-00}. ARC enables a single credential issuance to support multiple unlinkable presentations, up to a configurable limit. Each presentation uses a fresh nonce and randomization to ensure unlinkability, while the server enforces rate limits by tracking presentation-specific tags. 

Since SPRUCE provides a signature scheme with efficient proof protocols, we aim to construct a post-quantum SEP and subsequently a post-quantum variant of ARC.

@Carsten: not sure about the proofs. I added more context on ARC in~\Cref{sec:anonymous-credentials}
}